\documentclass[../../../include/open-logic-section]{subfiles}

\begin{document}
\olfileid[fr]{sth}{card-arithmetic}{simp}

\olsection{Simplification de l'addition et de la multiplication}

Il s'avère que l'addition et la multiplication
cardinales transfinies sont \emph{extrêmement}
simples. Cela tient au fait que les cardinaux sont
des ordinaux (d'un certain type), donc bien ordonnés,
et peuvent ainsi être manipulés d'une certaine manière.
Le démontrer, en revanche, n'est \emph{pas} si simple.
Commençons par une définition subtile :

\begin{defn}
Nous définissons un \emph{ordre canonique},
$\canonord$, sur les couples d'ordinaux en stipulant
que $\tuple{\alpha_1, \alpha_2} \canonord
\tuple{\beta_1, \beta_2}$ si et seulement si
l'une des conditions suivantes est remplie :
\begin{enumerate}
	\item $\max(\alpha_1, \alpha_2) < \max(\beta_1, \beta_2)$ ; ou
	\item $\max(\alpha_1, \alpha_2) = \max(\beta_1, \beta_2)$
	et $\alpha_1 < \beta_1$ ; ou
	\item $\max(\alpha_1, \alpha_2) = \max(\beta_1, \beta_2)$
	et $\alpha_1 = \beta_1$ et $\alpha_2 < \beta_2$.
\end{enumerate}
\end{defn}

\begin{lem}
$\tuple{\alpha \times \alpha, \canonord}$ est un
bon ordre pour tout ordinal $\alpha$.
\end{lem}

\begin{proof}
La relation $\canonord$ est manifestement connexe
sur $\alpha \times \alpha$. En effet, supposons
que ni $\tuple{\alpha_1, \alpha_2}$ ni
$\tuple{\beta_1, \beta_2}$ ne précède l'autre
selon $\canonord$. Alors
$\max(\alpha_1, \alpha_2) = \max(\beta_1, \beta_2)$,
$\alpha_1 = \beta_1$ et
$\alpha_2 = \beta_2$ ; donc
$\tuple{\alpha_1, \alpha_2} =
\tuple{\beta_1, \beta_2}$.

Pour établir que c'est un bon ordre, soit
$X \subseteq \alpha\times\alpha$ non vide.
Puisque $\alpha$ est un ordinal, il existe un plus
petit élément $\delta$ de
$\Setabs{\max(\gamma_1, \gamma_2)}
{\tuple{\gamma_1, \gamma_2} \in X}$.
Parmi les couples de
$\Setabs{\tuple{\gamma_1,\gamma_2} \in X}
{\max(\gamma_1, \gamma_2) = \delta}$,
ne gardons que ceux dont la première coordonnée
est minimale ; parmi eux, le couple dont la seconde
coordonnée est minimale est le plus petit élément
de $X$ pour $\canonord$.
\end{proof}
\noindent
Voici maintenant une observation toute simple :

\begin{prop}\ollabel{simplecardproduct}
Si $\cardeq{\alpha}{\beta}$, alors
$\cardeq{\alpha \times \alpha}{\beta \times \beta}$.
\end{prop}

\begin{proof}
Toute bijection
$f \colon \alpha \to \beta$ induit la bijection
$\tuple{\gamma_1, \gamma_2} \mapsto
\tuple{f(\gamma_1), f(\gamma_2)}$.
\end{proof}
\noindent
Nous pouvons désormais utiliser tout cela pour
démontrer un lemme essentiel :

\begin{lem}\ollabel{alphatimesalpha}
$\cardeq{\alpha}{\alpha \times \alpha}$,
pour tout ordinal infini $\alpha$.
\end{lem}

\begin{proof}
Raisonnons par l'absurde : soit $\alpha$ le plus
petit ordinal infini pour lequel l'énoncé est faux.
La \olref[sfr][siz][zigzag]{natsquaredenumerable}
montre que
$\cardeq{\omega}{\omega\times\omega}$ ;
donc $\omega \in \alpha$.
De plus, $\alpha$ est un cardinal. Sinon,
$\card{\alpha} \in \alpha$ ; l'hypothèse de
minimalité donnerait
$\cardeq{\card{\alpha}}
{\card{\alpha} \times \card{\alpha}}$.
Or $\cardeq{\card{\alpha}}{\alpha}$ par définition ;
donc la \olref{simplecardproduct} entraînerait
$\cardeq{\alpha}{\alpha\times\alpha}$,
ce qui est contradictoire.

Pour chaque
$\tuple{\gamma_1, \gamma_2} \in \alpha \times \alpha$,
considérons le segment :
\begin{align*}
	\text{Seg}(\gamma_1, \gamma_2)
	&= \Setabs{\tuple{\delta_1, \delta_2}
	\in \alpha \times \alpha}
	{\tuple{\delta_1, \delta_2}
	\canonord \tuple{\gamma_1, \gamma_2}}
\end{align*}
Posons $\gamma = \max(\gamma_1, \gamma_2)$.
Alors $\tuple{\gamma_1, \gamma_2}
\canonord \tuple{\gamma+1, \gamma + 1}$.
Si $\gamma$ est infini, on obtient :
\begin{align*}
	\text{Seg}(\gamma_1, \gamma_2) &
	\precsim ((\gamma \ordplus 1)
	\ordtimes (\gamma \ordplus 1))\\
	&\approx (\gamma \ordtimes \gamma)
	\text{, par le
	\olref[ord-arithmetic][using-addition]{ordinfinitycharacter}
	et la \olref{simplecardproduct}}\\
	&\approx \gamma
	\text{, par l'hypothèse de minimalité}\\
	& \prec \alpha
	\text{, puisque $\alpha$ est un cardinal}
\end{align*}



Donc $\ordtype{\alpha\times \alpha, \canonord}
\leq \alpha$, et par conséquent
$\cardle{\alpha \times \alpha}{\alpha}$.
Comme, bien entendu,
$\cardle{\alpha}{\alpha \times \alpha}$,
le résultat découle du théorème de
Schröder--Bernstein.
\end{proof}

Nous obtenons enfin le résultat qui simplifie
les opérations :

\begin{thm}\ollabel{cardplustimesmax}
Si $\cardfont{a}, \cardfont{b}$ sont des
cardinaux infinis, alors :
\[
	\cardfont{a}
\cardtimes \cardfont{b}
= \cardfont{a} \cardplus \cardfont{b}
= \text{max}(\cardfont{a}, \cardfont{b}).
\]
\end{thm}

\begin{proof}
Sans perte de généralité, supposons que
$\cardfont{a} = \max(\cardfont{a}, \cardfont{b})$.
Alors, par le \olref{alphatimesalpha},
$\cardfont{a}\cardtimes\cardfont{a}
= \cardfont{a}
\leq \cardfont{a}\cardplus \cardfont{b}
\leq \cardfont{a} \cardplus \cardfont{a}
\leq \cardfont{a} \cardtimes\cardfont{a}$.
\end{proof}
\noindent
De même, si $\cardfont{a}$ est infini,
l'union de $\cardfont{a}$ ensembles ayant
chacun au plus $\cardfont{a}$ éléments
possède elle aussi au plus $\cardfont{a}$ éléments :

\begin{prop}\ollabel{kappaunionkappasize}
Soit $\cardfont{a}$ un cardinal infini.
Pour chaque ordinal $\beta \in \cardfont{a}$,
soit $X_\beta$ un ensemble tel que
$\card{X_\beta} \leq \cardfont{a}$.
Alors
$\card{\bigcup_{\beta \in \cardfont{a}} X_\beta}
\leq \cardfont{a}$.
\end{prop}

\begin{proof}
Pour chaque $\beta \in \cardfont{a}$,
fixons !!a{injection}
$f_\beta \colon X_\beta \to \cardfont{a}$.\footnote{
Comment fixer toutes ces injections ?
Voir la \olref[sth][choice][countablechoice]{sec}.}
Définissons !!a{injection}
$g \colon \bigcup_{\beta \in \cardfont{a}} X_\beta
\to \cardfont{a} \times \cardfont{a}$
par $g(v) = \tuple{\beta, f_\beta(v)}$,
où $v \in X_\beta$ et où
$v \notin X_\gamma$ pour tout $\gamma \in \beta$.
Alors
$\bigcup_{\beta \in \cardfont{a}} X_\beta
\preceq \cardfont{a} \times \cardfont{a}
\approx \cardfont{a}$,
par le \olref{cardplustimesmax}.
\end{proof}

\end{document}
